% ============================================================================
% 与项目其他模块的关系
% ============================================================================

\subsection{从底层到大规模并行}

这个项目的各个模块构成了从最底层硬件到大规模并行计算的完整学习路径：

\begin{enumerate}
    \item \textbf{Chip（Verilog 硬件设计）}：最底层，理解数字电路如何工作
    \item \textbf{ISA（AArch64 汇编）}：理解 CPU 指令集和单线程程序执行
    \item \textbf{QEMU（裸机编程）}：理解没有操作系统时程序如何直接与硬件交互
    \item \textbf{Algo（C 语言算法）}：用高级语言实现经典算法，理解算法复杂度
    \item \textbf{CUDA（并行计算）}：理解如何利用成千上万个核心并行解决问题
\end{enumerate}

\subsection{CUDA 与汇编的对比}

\begin{itemize}
    \item \textbf{ISA（汇编）}：关注单线程的精确控制，每一条指令都清楚 CPU 在做什么
    \item \textbf{CUDA}：关注大规模并行，把问题分解成几千个线程同时执行
\end{itemize}

两者互补：汇编让你理解单个计算单元的细节，CUDA 让你理解如何组织大量计算单元协作。

\subsection{与算法的关系}

第 5 章的算法都是串行的，而很多算法天然可以并行化：
\begin{itemize}
    \item 排序（基数排序、快速排序的某些阶段）
    \item 矩阵运算
    \item 图遍历（BFS）
    \item 动态规划的某些问题
\end{itemize}

理解了串行算法后，再学习如何用 CUDA 加速它们是很自然的进阶路径。
